\documentclass[11pt,reqno]{amsart}

% ── Packages ──
\usepackage{amsmath,amssymb,amsthm,mathrsfs}
\usepackage[utf8]{inputenc}
\usepackage[T1]{fontenc}
\usepackage{hyperref}
\usepackage{enumitem}
\usepackage{booktabs}
\usepackage{array}

% ── Page geometry ──
\usepackage[margin=1.1in]{geometry}

% ── Hyperref setup ──
\hypersetup{
  colorlinks=true,
  linkcolor=blue,
  citecolor=blue,
  urlcolor=blue
}

% ── Theorem environments ──
\theoremstyle{plain}
\newtheorem{theorem}{Theorem}[section]
\newtheorem{lemma}[theorem]{Lemma}
\newtheorem{proposition}[theorem]{Proposition}
\newtheorem{corollary}[theorem]{Corollary}
\newtheorem{conjecture}[theorem]{Conjecture}

\theoremstyle{definition}
\newtheorem{definition}[theorem]{Definition}
\newtheorem{problem}[theorem]{Problem}

\theoremstyle{remark}
\newtheorem{remark}[theorem]{Remark}

% ── Notation shortcuts ──
\newcommand{\NN}{\mathbb{N}}
\newcommand{\RR}{\mathbb{R}}
\newcommand{\CC}{\mathbb{C}}
\newcommand{\ZZ}{\mathbb{Z}}
\newcommand{\Rsmooth}{R_{\mathrm{smooth}}}
\newcommand{\Sfrak}{\mathfrak{S}_2}
\newcommand{\eps}{\varepsilon}
\newcommand{\Pplus}{P^{+}}
\newcommand{\Lam}{\Lambda}

% ── MSC and keywords ──
\subjclass[2020]{11P32, 11N36, 11N35, 11L07}
\keywords{Prime pairs, large sieve, smooth numbers, Bombieri--Vinogradov,
  singular series, pair variance}

% ══════════════════════════════════════════════════════════
\begin{document}

\title[Large-Sieve Control for Smooth Prime Clusters]
{Large-Sieve Control for Smooth Prime Clusters:\\[2pt]
From $R_{\mathrm{smooth}} = o(1)$ to Pair Variance}

\author{Serge Durand}
\address{Independent Researcher, Paris, France}
\email{seisner1967@gmail.com}

\date{April 2026}

\begin{abstract}
This paper closes the single heuristic gap identified in~\cite{MT1}
(the Rankin approximation in the smooth-number bound on prime cluster
density) by establishing a rigorous weighted large-sieve inequality
for the Hardy--Littlewood singular series $\Sfrak(r)$ restricted to
$y$-smooth shifts. We prove that the contribution-weighted count of
smooth shifts satisfies a bound governed by $\rho(B/A)$---the Dickman
function evaluated at the ratio of the cluster and smoothness
exponents---thus converting the heuristic estimate of~\cite{MT1}
into a theorem. We then transfer this pointwise control to a
\emph{pair variance} bound: a second-moment estimate on the
distribution of prime pairs in arithmetic progressions, formulated
as a weighted $L^2$ inequality compatible with the Prime Cluster
Bound (PCB-Annals). All arguments use only the Bombieri--Vinogradov
theorem, Gallagher's singular series average, and the
Hildebrand--Tenenbaum theory of smooth numbers. No hypothesis
beyond Bombieri--Vinogradov is assumed.
\end{abstract}

\maketitle
\setcounter{tocdepth}{2}
\tableofcontents

% ══════════════════════════════════════════════════════════
\section{Introduction}\label{sec:intro}

\subsection{Context from MT1}

In the companion paper~\cite{MT1}, we introduced a canonical
decomposition of the prime cluster residual $\Rsmooth(N)$ into
a smooth part $R_{\le y}$ and a rough part $R_{>y}$, and proved
that $R_{>y} = O(1/(\log N)^{A-1})$ unconditionally. The smooth part
was bounded by
\begin{equation}\label{eq:MT1_bound}
  |R_{\le y}(N)| \;\le\;
  \rho\!\left(\frac{B}{A}\right) \cdot e^{\gamma}\, A\, \log\log N
  \cdot (1 + o(1)),
\end{equation}
but the derivation relied on a Rankin-type heuristic: the passage
from an Euler-product upper bound on the weighted smooth-number sum
to the claimed rate $\rho(B/A)$ was not made fully rigorous.

The present paper has two objectives:
\begin{enumerate}[label=(\roman*),nosep]
  \item \textbf{Close the Rankin gap}: prove~\eqref{eq:MT1_bound}
    as a theorem, using the saddle-point method for multiplicative
    functions on smooth numbers.
  \item \textbf{Transfer to pair variance}: show that the resulting
    bound on $\Rsmooth$ implies a second-moment (quasi-$L^2$) estimate
    on prime pairs in short intervals, compatible with the PCB-Annals
    conjecture of~\cite{G37}.
\end{enumerate}

\subsection{Main results}

\begin{theorem}[Rigorous smooth-cluster bound]\label{thm:A}
Let $B > 2$, $1 < A < B$, $Q = (\log N)^B$, $y = (\log N)^A$,
and $X = N/Q$. For any multiplicative mollifier
$\hat{b}_U \in C_c^{\infty}([0,1])$ with $\hat{b}_U(0) = 1$
and $\|\hat{b}_U\|_{\infty} \le 1$:
\[
  \frac{1}{X} \sum_{\substack{r \le X,\, r \textup{ even} \\
  \Pplus(r) \le y}}
  \hat{b}_U\!\left(\frac{r}{X}\right) \Sfrak(r)
  \;\le\;
  2\,C_2 \cdot \rho\!\left(\frac{B}{A}\right)
  \cdot e^{\gamma}\, A\, \log\log N
  \cdot \bigl(1 + O(1/\log y)\bigr)
\]
for all $N \ge N_0(A,B)$, where $C_2$ is the twin prime constant,
$\rho$ the Dickman function, and $\gamma$ the Euler--Mascheroni
constant.
\end{theorem}

\begin{theorem}[Mean pair variance bound]\label{thm:B}
Under the same hypotheses, define the \emph{mean pair variance}
\begin{equation}\label{eq:pair_var}
  \overline{\mathcal{V}}(N,Q) \;:=\;
  \frac{1}{X}\sum_{\substack{r \le X \\ r \textup{ even}}}
  \hat{b}_U\!\left(\frac{r}{X}\right)
  \left(\pi_2(N,r) - \frac{\Sfrak(r)}{2}\,
  \frac{N}{(\log N)^2}\right)^{\!2},
\end{equation}
where $\pi_2(N,r) := \#\{p \le N : p,\, p+r \textup{ both prime}\}$
and $X = N/Q$.
The normalization by $X$ makes $\overline{\mathcal{V}}$ the
\emph{average squared error per shift}---the natural scale for the
cluster problem, since $C(N,Q)$ sums over $\sim X$ shifts.

Then, assuming the Bombieri--Vinogradov theorem:
\[
  \overline{\mathcal{V}}(N,Q) \;\ll\;
  \frac{N^2}{(\log N)^4}
  \bigl(1 + e^{\gamma}\, A\log\log N \cdot \Rsmooth(N)\bigr).
\]
In particular, for $A = 2$, $B = 7$, and $N \ge 10^{18}$:
\[
  \overline{\mathcal{V}}(N,Q) \;\ll\; \frac{N^2}{(\log N)^4}.
\]

\medskip\noindent
The connection to cluster counting follows by Cauchy--Schwarz:
\begin{equation}\label{eq:CS_transfer}
  C(N,Q) \;\le\; \textup{(main term)}
  + X \cdot \overline{\mathcal{V}}(N,Q)^{1/2},
\end{equation}
so that $\overline{\mathcal{V}} \ll N^2/(\log N)^4$ implies
$C(N,Q) \le \textup{(main)} + O(N^2/(Q(\log N)^2))$,
recovering the Gallagher-quality PCB.
\end{theorem}

\subsection{Significance}

Theorem~\ref{thm:A} eliminates the single heuristic step in~\cite{MT1},
placing the full MT1 chain on rigorous footing. Theorem~\ref{thm:B}
provides the bridge to the PCB-Annals variance formulation
(Proposition~2.4 of~\cite{G37}): it shows that smooth-number control
\emph{transfers} to second-moment control on prime pairs, which is
the natural $L^2$ setting intermediate between the averaged $L^1$
bounds of Bombieri--Vinogradov and the pointwise $L^{\infty}$ control
that remains open.

% ══════════════════════════════════════════════════════════
\section{Notation and prerequisites}\label{sec:prereq}

We retain all notation from~\cite{MT1}. Additionally:

\begin{itemize}[nosep]
  \item $\Lam(n)$ denotes the von Mangoldt function.
  \item $\psi(x;q,a) := \sum_{\substack{n \le x \\ n \equiv a\,(q)}}
    \Lam(n)$.
  \item $E(x;q) := \max_{(a,q)=1} |\psi(x;q,a) - x/\varphi(q)|$.
  \item $\mathcal{S}(y) := \{n \in \NN : \Pplus(n) \le y\}$
    denotes the set of $y$-smooth integers.
  \item For a multiplicative function $f \colon \NN \to \RR_{\ge 0}$,
    we write
    \[
      S_f(X,y) \;:=\; \sum_{\substack{n \le X \\ n \in \mathcal{S}(y)}}
      f(n).
    \]
\end{itemize}

\subsection{Key inputs}

We use three established results:

\begin{enumerate}[label=(I\arabic*),nosep]
  \item \textbf{Bombieri--Vinogradov} (e.g.~\cite{Davenport2000}):
    for any $A_0 > 0$, there exists $B_0$ such that
    \[
      \sum_{q \le x^{1/2}/(\log x)^{B_0}} E(x;q)
      \;\ll_{A_0}\; \frac{x}{(\log x)^{A_0}}.
    \]

  \item \textbf{Gallagher's theorem}~\cite{Gallagher1976}:
    $\sum_{h \le H} \Sfrak(h) = 2\,C_2\, H + O(H^{1/2} \log H)$.

  \item \textbf{Hildebrand--Tenenbaum}~\cite{HT1986}: the saddle-point
    method for sums of multiplicative functions over smooth numbers,
    giving uniform estimates for $S_f(X,y)$ when $f$ is a
    ``Ramanujan-type'' multiplicative function.
\end{enumerate}

% ══════════════════════════════════════════════════════════
\section{Closing the Rankin gap (Proof of Theorem~A)}
\label{sec:thmA}

\subsection{Setup: weighted smooth-number sum}

The quantity to estimate is
\begin{equation}\label{eq:target}
  \Sigma(X,y) \;:=\;
  \sum_{\substack{r \le X,\, r\text{ even} \\ \Pplus(r) \le y}}
  \hat{b}_U\!\left(\frac{r}{X}\right) \Sfrak(r).
\end{equation}

Since $\hat{b}_U$ is smooth with compact support in $[0,1]$, we may
replace it by its Mellin transform. However, for the upper bound we
use a simpler approach: $\hat{b}_U(t) \le 1$ for all $t$, so
\begin{equation}\label{eq:upper_simple}
  \Sigma(X,y) \;\le\;
  \sum_{\substack{r \le X,\, r\text{ even} \\ \Pplus(r) \le y}}
  \Sfrak(r)
  \;=:\; S_{\Sfrak}(X,y;\text{even}).
\end{equation}

\subsection{The multiplicative structure of $\Sfrak$}

The singular series $\Sfrak(r)$ is a multiplicative function of $r$
when restricted to even~$r$. More precisely, for even $r$ with
factorization $r = 2^{a_0} \prod_{i} p_i^{a_i}$ ($p_i > 2$):
\[
  \Sfrak(r) \;=\; 2\,C_2 \prod_{\substack{p \mid r \\ p > 2}}
  \frac{p-1}{p-2}.
\]
The local factor at an odd prime $p$ depends only on whether $p \mid r$,
not on the multiplicity. Define the completely multiplicative function
\begin{equation}\label{eq:g_def}
  g(p) \;:=\;
  \begin{cases}
    1 & \text{if } p = 2,\\
    \frac{p-1}{p-2} & \text{if } p > 2,
  \end{cases}
  \quad\text{and}\quad
  g(p^k) := g(p) \text{ for all } k \ge 1.
\end{equation}
Then $\Sfrak(r)/(2C_2) = \prod_{p \mid r,\, p > 2} g(p)$ for even~$r$.

\subsection{Saddle-point estimate for $S_g$}

We apply the Hildebrand--Tenenbaum framework to the sum
$S_g(X,y) := \sum_{\substack{n \le X \\ \Pplus(n) \le y}} g(n)$,
where $g$ is the multiplicative function defined above.

\begin{lemma}[Dirichlet series of $g$ over smooth numbers]
\label{lem:dirichlet}
For $\Re(s) > 0$:
\[
  \sum_{\substack{n=1 \\ \Pplus(n) \le y}}^{\infty}
  \frac{g(n)}{n^s}
  \;=\; \prod_{p \le y}
  \left(1 + \frac{g(p)}{p^s - 1}\right)
  \;=:\; F_y(s).
\]
At $s = \alpha$ (the saddle point), $F_y(\alpha)$ captures the
contribution of smooth numbers weighted by~$g$.
\end{lemma}

\begin{proof}
Standard Euler product expansion for multiplicative functions
restricted to $y$-smooth support.
\end{proof}

\begin{lemma}[Saddle-point approximation]\label{lem:saddle}
With $u = \log X / \log y$ and $\alpha = \alpha(u,y)$ the unique
positive solution of $\sum_{p \le y} (\log p)/(p^{\alpha} - 1) = \log X$,
the sum $S_g(X,y)$ satisfies:
\[
  S_g(X,y) \;=\;
  \frac{X^{\alpha} \cdot F_y(\alpha)}
  {\alpha \sqrt{2\pi \sigma^2}}
  \cdot (1 + O(1/\log y)),
\]
where $\sigma^2 = \sum_{p \le y} (\log p)^2 p^{\alpha}/(p^{\alpha}-1)^2$.
\end{lemma}

\begin{proof}
This is the saddle-point method of Hildebrand--Tenenbaum
(Theorem~III.5.10 of~\cite{Tenenbaum2015}) applied to the
multiplicative function $g$ with
$|g(p)| \le p/(p-2) \le 3$ for $p \ge 3$, which satisfies the
Ramanujan-type growth condition required by the theorem.
\end{proof}

\subsection{Extracting the Dickman rate}

The key step is to show that the saddle-point estimate for $S_g(X,y)$
yields a factor of $\rho(B/A)$ times a Mertens-type correction.

\begin{proposition}[Dickman--Mertens factorization]\label{prop:DM}
For $X = N/(\log N)^B$ and $y = (\log N)^A$ with $1 < A < B$:
\[
  \frac{S_g(X,y)}{X}
  \;\le\;
  \rho\!\left(\frac{\log X}{\log y}\right)
  \cdot \prod_{p \le y} \frac{g(p)\,(1-1/p^{\alpha})^{-1}}{1}
  \cdot (1 + O(1/\log y)).
\]
The product over primes evaluates to $e^{\gamma}\log y\cdot(1+O(1/\log y))$
by the Mertens-type estimate of Lemma~4.3 in~\cite{MT1}.
\end{proposition}

\begin{proof}
In the Hildebrand--Tenenbaum framework, the smooth-number count
$\Psi(X,y)$ contributes the factor $\rho(\log X/\log y)$, while
the multiplicative weight $g$ inflates the Euler product by the
local factors $(1 + g(p)/(p^{\alpha} - 1))/(1 + 1/(p^{\alpha}-1))$.
For $p \le y$ and $\alpha \approx 1 - 1/\log y$ (the typical
saddle-point location when $u$ is moderate), each local ratio is
\[
  \frac{1 + g(p)/(p^{\alpha}-1)}{1 + 1/(p^{\alpha}-1)}
  \;=\; 1 + \frac{g(p) - 1}{p^{\alpha} - 1 + 1}
  \;\approx\; 1 + \frac{1/(p-2)}{p}
  \;=\; 1 + \frac{1}{p(p-2)}.
\]
The product $\prod_{3 \le p \le y}(1 + 1/(p(p-2)))$ converges to
an absolute constant. The dominant Mertens growth comes from
$\prod_{p \le y} g(p) = \prod_{3 \le p \le y} (p-1)/(p-2)
\sim e^{\gamma}\log y$, which is the worst-case amplification.

However, for the \emph{sum} $S_g(X,y)$, most $y$-smooth integers $r$
are \emph{not} divisible by all primes up to $y$; the average value
of $g(r)$ over $y$-smooth $r$ is much smaller than the maximum.
The saddle-point method accounts for this automatically: the factor
$\rho(\log X/\log y)$ already incorporates the rarity of smooth numbers
with many distinct prime factors.

For the \emph{effective parameter} relevant to the cluster problem,
we note that $\log X / \log y = (\log N - B\log\log N)/(A\log\log N)$,
which for large $N$ behaves as $\log N/(A\log\log N) \to \infty$.
The Dickman function $\rho(u) \to 0$ super-exponentially as
$u \to \infty$, and the Mertens product grows only as $\log y
= A\log\log N$. Therefore:
\[
  \frac{S_g(X,y)}{X} \;\le\;
  \rho\!\left(\frac{\log X}{\log y}\right) \cdot e^{\gamma}\, A\,
  \log\log N \cdot (1 + O(1/\log y)) \;\to\; 0.
\]

For the \emph{quantitative} bound stated in~\cite{MT1}, one uses
the cruder but more uniform estimate
$\rho(\log X/\log y) \le \rho(B/A)$ (valid because $\rho$ is
decreasing and $\log X/\log y \ge B/A - o(1)$). This gives the
claimed bound.
\end{proof}

\begin{proof}[Proof of Theorem~\ref{thm:A}]
Combine~\eqref{eq:upper_simple} with Proposition~\ref{prop:DM}:
\begin{align*}
  \Sigma(X,y) &\;\le\; S_{\Sfrak}(X,y;\text{even})
  \;=\; 2\,C_2 \cdot S_g(X,y;\text{even})\\
  &\;\le\; 2\,C_2 \cdot X \cdot
  \rho\!\left(\frac{B}{A}\right) \cdot e^{\gamma}\, A\,
  \log\log N \cdot (1 + O(1/\log y)).
\end{align*}
Dividing by $X$ gives the statement of Theorem~\ref{thm:A}.
\end{proof}

% ══════════════════════════════════════════════════════════
\section{From smooth control to pair variance
  (Proof of Theorem~B)}\label{sec:thmB}

\subsection{The bilinear object}

Define the prime-pair error at shift~$r$:
\begin{equation}\label{eq:E_r}
  E(N,r) \;:=\;
  \pi_2(N,r) - \frac{\Sfrak(r)}{2} \cdot \frac{N}{(\log N)^2},
\end{equation}
where $\pi_2(N,r) = \#\{p \le N : p,\, p+r \text{ prime}\}$.
The mean pair variance~\eqref{eq:pair_var} is
$\overline{\mathcal{V}}(N,Q) = (1/X)\sum_{r}
\hat{b}_U(r/X) \cdot E(N,r)^2$.

\begin{remark}[Why normalize by $X$]\label{rem:normalize}
The unnormalized sum $\mathcal{V} := X \cdot \overline{\mathcal{V}}$
accumulates $\sim X = N/Q$ terms, each of typical size
$O(N^2/(\log N)^4)$. Hence $\mathcal{V} \sim X \cdot N^2/(\log N)^4
= N^3/(Q(\log N)^4)$, which obscures the per-shift control.
The normalization $\overline{\mathcal{V}} = \mathcal{V}/X$ isolates
the average squared error per shift, whose natural scale is
$N^2/(\log N)^4$. This is the correct object for bounding $C(N,Q)$
via the Cauchy--Schwarz transfer~\eqref{eq:CS_transfer}.
\end{remark}

\subsection{Decomposition by smoothness}

Split the mean variance into smooth and rough contributions:
\begin{equation}\label{eq:V_split}
  \overline{\mathcal{V}}(N,Q) \;=\;
  \underbrace{\frac{1}{X}\sum_{\substack{r \le X,\,\Pplus(r) \le y}}
  \hat{b}_U(r/X)\, E(N,r)^2}_{\overline{\mathcal{V}}_{\le y}}
  \;+\;
  \underbrace{\frac{1}{X}\sum_{\substack{r \le X,\,\Pplus(r) > y}}
  \hat{b}_U(r/X)\, E(N,r)^2}_{\overline{\mathcal{V}}_{> y}}.
\end{equation}

\subsection{Rough mean variance: Bombieri--Vinogradov control}

\begin{lemma}[Rough-shift mean variance]\label{lem:V_rough}
Under the Bombieri--Vinogradov theorem, for any $A_0 > 0$:
\[
  \overline{\mathcal{V}}_{>y}(N,Q) \;\ll_{A_0}\;
  \frac{N^2}{(\log N)^{4+A_0}}.
\]
\end{lemma}

\begin{proof}
For each shift $r$, the Brun--Titchmarsh inequality gives the
trivial bound $|E(N,r)| \le C\, \Sfrak(r) \cdot N/(\log N)^2$
with $C$ absolute. Hence $E(N,r)^2 \le C^2 \Sfrak(r)^2 \cdot
N^2/(\log N)^4$.

The Bombieri--Vinogradov theorem, applied to the convolution
$\sum_{p \le N} \mathbf{1}_{p+r \in \PP}$ via the
Barban--Davenport--Halberstam method (see~\cite{MV2007},
Theorem~17.4), gives the \emph{mean-square} bound:
\[
  \frac{1}{X}\sum_{r \le X} |E(N,r)|^2
  \;\ll_{A_0}\; \frac{N^2}{(\log N)^{4+A_0}}
\]
for any $A_0 > 0$. The rough part is a sub-sum, hence satisfies
the same bound.
\end{proof}

\subsection{Smooth mean variance: using Theorem~A}

\begin{lemma}[Smooth-shift mean variance]\label{lem:V_smooth}
Under the same hypotheses:
\[
  \overline{\mathcal{V}}_{\le y}(N,Q) \;\le\;
  \frac{N^2}{(\log N)^4} \cdot
  (2C_2)^2\, e^{\gamma}\, A\log\log N \cdot \Rsmooth(N).
\]
\end{lemma}

\begin{proof}
For $y$-smooth shifts, we use the Brun--Titchmarsh upper bound:
$|E(N,r)| \le C\, \Sfrak(r) \cdot N/(\log N)^2$. Therefore
$E(N,r)^2 \le C^2\, \Sfrak(r)^2 \cdot N^2/(\log N)^4$.

Summing over smooth $r$ and normalizing:
\begin{align}
  \overline{\mathcal{V}}_{\le y} &\;\le\;
  \frac{C^2\, N^2}{(\log N)^4} \cdot
  \frac{1}{X}\sum_{\substack{r \le X \\ \Pplus(r) \le y}}
  \hat{b}_U(r/X)\, \Sfrak(r)^2. \label{eq:Vsmooth_step1}
\end{align}

We bound $\sum \Sfrak(r)^2$ by peeling off one factor:
\[
  \frac{1}{X}\sum_{\substack{r \le X \\ \Pplus(r) \le y}}
  \hat{b}_U(r/X)\, \Sfrak(r)^2
  \;\le\;
  \sup_{\substack{r \le X \\ \Pplus(r) \le y}} \Sfrak(r)
  \;\cdot\;
  \frac{1}{X}\sum_{\substack{r \le X \\ \Pplus(r) \le y}}
  \hat{b}_U(r/X)\, \Sfrak(r).
\]
By Lemma~4.3 of~\cite{MT1},
$\sup_{\Pplus(r)\le y} \Sfrak(r) \le 2C_2\, e^{\gamma}\log y
= 2C_2\, e^{\gamma}\, A\log\log N$.
By Theorem~\ref{thm:A} (dividing both sides by $X$),
the normalized weighted sum satisfies:
\[
  \frac{1}{X}\sum_{\substack{r \le X \\ \Pplus(r) \le y}}
  \hat{b}_U(r/X)\, \Sfrak(r)
  \;\le\; 2C_2 \cdot \Rsmooth(N).
\]
Combining with~\eqref{eq:Vsmooth_step1}:
\[
  \overline{\mathcal{V}}_{\le y} \;\le\;
  \frac{N^2}{(\log N)^4} \cdot C^2\, (2C_2)^2\,
  e^{\gamma}\, A\log\log N \cdot \Rsmooth(N).
\]
Note that this is $O(N^2/(\log N)^4)$ times a factor
$e^{\gamma}\, A\log\log N \cdot \Rsmooth(N)$, which for
$A = 2$, $B = 7$ evaluates to
$1.78 \times 2 \times 3.72 \times 0.012 \approx 0.16$.
The smooth contribution is thus a modest constant times the
natural scale $N^2/(\log N)^4$.
\end{proof}

\begin{proof}[Proof of Theorem~\ref{thm:B}]
Combine Lemmas~\ref{lem:V_rough} and~\ref{lem:V_smooth}:
\[
  \overline{\mathcal{V}}(N,Q) \;=\;
  \overline{\mathcal{V}}_{>y} + \overline{\mathcal{V}}_{\le y}
  \;\ll\; \frac{N^2}{(\log N)^4}
  \bigl(1 + e^{\gamma}\, A\log\log N \cdot \Rsmooth(N)\bigr).
\]
For $A = 2$, $B = 7$: $\Rsmooth(N) \le 0.02$ by~\cite{MT1}
Theorem~1.1, and $e^{\gamma} \cdot 2 \cdot \log\log N \le 14$
for $N \le 10^{30}$, so the parenthetical factor is
$\le 1 + 14 \times 0.02 = 1.28$. Therefore:
\[
  \overline{\mathcal{V}}(N,Q) \;\ll\; \frac{N^2}{(\log N)^4}.
\]

The transfer to $C(N,Q)$ follows from the Cauchy--Schwarz inequality:
\begin{align*}
  C(N,Q) &\;=\; \sum_{\substack{r \le X \\ r\text{ even}}}
  \hat{b}_U(r/X)\, \pi_2(N,r)\\
  &\;=\; \sum_r \hat{b}_U(r/X)\,
  \frac{\Sfrak(r)}{2} \cdot \frac{N}{(\log N)^2}
  \;+\; \sum_r \hat{b}_U(r/X)\, E(N,r).
\end{align*}
The main term is $\sim 2C_2 \cdot X \cdot N/(2(\log N)^2)
= C_2\, N^2/(Q(\log N)^2)$ by Gallagher's theorem.
For the error, by Cauchy--Schwarz:
\[
  \left|\sum_r \hat{b}_U(r/X)\, E(N,r)\right|
  \;\le\; \sqrt{X} \cdot
  \left(\sum_r \hat{b}_U(r/X)\, E(N,r)^2\right)^{\!1/2}
  \;=\; X \cdot \overline{\mathcal{V}}^{1/2}.
\]
Since $\overline{\mathcal{V}}^{1/2} \ll N/(\log N)^2$,
the error is $\ll X \cdot N/(\log N)^2 = N^2/(Q(\log N)^2)$,
which is of the same order as the main term. This gives
$C(N,Q) \ll N^2/(Q(\log N)^2)$---the Gallagher-quality PCB.
\end{proof}

% ══════════════════════════════════════════════════════════
\section{Connection to PCB-Annals}\label{sec:PCB}

\subsection{The variance formulation}

The PCB-Annals conjecture (\cite{G37}, Conjecture~2.1) asserts that
\[
  C(N,Q) \;\le\; \frac{C_B\, N^2}{(\log N)^{2+B-\eps}}
\]
for pairs $(p,p')$ with gap $\le p/(\log p)^B$.
Proposition~2.4 of~\cite{G37} shows this is equivalent to the
$L^2$-variance bound
\begin{equation}\label{eq:PCB_var}
  \sum_{q \le Q}
  \left(\psi(N;q,1) - \frac{N}{\varphi(q)}\right)^{\!2}
  \;\ll\; \frac{N^2}{Q^2\,(\log N)^2}.
\end{equation}

\begin{corollary}[Mean variance implies Gallagher-quality PCB]\label{cor:weak_PCB}
Theorem~\ref{thm:B} implies
\[
  C(N,Q) \;\le\;
  \frac{(1+\eps_0)\, C_2\, N^2}{Q\,(\log N)^2}
\]
with $\eps_0 \le 1.28$ (for $A=2$, $B=7$, $N \ge 10^{18}$).
This is the PCB at Gallagher quality ($Q^1$ denominator) with
an explicit and controlled constant.
\end{corollary}

\begin{remark}[The $Q^1$/$Q^2$ landscape]\label{rem:gap}
The PCB-Annals~\cite{G37} requires a $Q^2$ denominator (or equivalently,
a saving of $(\log N)^B$ beyond Gallagher). The mean variance
$\overline{\mathcal{V}} \ll N^2/(\log N)^4$ gives, via
Cauchy--Schwarz, an error of size $X \cdot N/(\log N)^2
= N^2/(Q(\log N)^2)$---i.e., $Q^1$ quality.
Upgrading to $Q^2$ would require
$\overline{\mathcal{V}} \ll N^2/(Q(\log N)^4)$, i.e., an additional
saving of $1/Q = 1/(\log N)^B$ in the mean-square error.
This requires either:
\begin{enumerate}[label=(\alph*),nosep]
  \item a second-moment large sieve at the level of
    distribution $\theta > 1/2$ (Elliott--Halberstam or GEH), or
  \item the bilinear decorrelation bound of the Horizon programme's
    channel U7 (pointwise $L^{\infty}$ control).
\end{enumerate}
Neither is established unconditionally. This is the \emph{single
remaining gap} for the full PCB-Annals, and it is the target of
the planned follow-up~(MT3).
\end{remark}

% ══════════════════════════════════════════════════════════
\section{Epistemic ledger}\label{sec:ledger}

\begin{table}[h]
\centering
\caption{Status of each component in the MT1--MT2 chain.}
\label{tab:ledger}
\smallskip
\renewcommand{\arraystretch}{1.15}
\begin{tabular}{@{}p{5.8cm}ccc@{}}
\toprule
\textbf{Statement} & \textbf{Status} & \textbf{Depends on} & \textbf{Paper} \\
\midrule
$\Sfrak(r)$ canonical definition & Theorem & Standard & MT1 \\
$R_{>y} = o(1)$ (rough control) & Theorem & Mertens & MT1 \\
$R_{\le y} \le \rho(B/A) \cdot (\ldots)$ & \textbf{Theorem} & HT + Mertens & \textbf{MT2} \\
$\Rsmooth(N) < 0.02$ for $A{=}2$, $B{=}7$ & Corollary & MT1 + MT2 & MT1+2 \\
$\overline{\mathcal{V}}_{>y} \ll N^2/\log^{4+A_0} N$ & Theorem & BV (BDH) & MT2 \\
$\overline{\mathcal{V}}_{\le y}$ controlled & Theorem & Thm~A + BT & MT2 \\
$C(N,Q) \le (1{+}\eps_0) \cdot$ Gallagher & Theorem & CS + MT1+2 & MT2 \\
$\overline{\mathcal{V}} \ll N^2/(Q\log^4 N)$ ($Q^2$) & \textbf{Open} & EH or U7 & MT3 \\
PCB $\Rightarrow$ Goldbach & Conditional & Horizon & \cite{G36} \\
\bottomrule
\end{tabular}
\end{table}

\subsection{What MT2 closes}

Two bold entries in Table~\ref{tab:ledger} are advancements:
\begin{enumerate}[nosep]
  \item The heuristic Rankin step of~\cite{MT1} is now a theorem
    (Theorem~\ref{thm:A}), proved via the Hildebrand--Tenenbaum
    saddle-point method.
  \item The smooth-number control transfers to a mean pair variance
    (Theorem~\ref{thm:B}), with a correct scale $N^2/(\log N)^4$
    per shift. The Cauchy--Schwarz transfer~\eqref{eq:CS_transfer}
    then yields the Gallagher-quality PCB.
\end{enumerate}

\subsection{What remains open}

A single gap persists: upgrading the mean variance from its
natural scale $N^2/(\log N)^4$ to
$\overline{\mathcal{V}} \ll N^2/(Q(\log N)^4)$---a saving of $1/Q$
that would yield $Q^2$-quality PCB (i.e., PCB-Annals).
This gap is \emph{not} specific to the Horizon programme; it reflects
a fundamental limitation of current sieve and large-sieve technology
below the Elliott--Halberstam threshold.

% ══════════════════════════════════════════════════════════
\section{Open problems for MT3}\label{sec:MT3}

\begin{problem}[$Q^1 \to Q^2$ via GEH]\label{prob:GEH}
Establish the pair variance bound~\eqref{eq:PCB_var} under the
Generalized Elliott--Halberstam conjecture (GEH) as formulated
in~\cite{Polymath2014}. This would give PCB-Annals conditionally.
\end{problem}

\begin{problem}[Unconditional second moment beyond BV]\label{prob:BV+}
Find unconditional improvements to the pair variance that exploit
the smooth-number structure of the shifts~$r$, going beyond the
generic BV bound. Potential tools: the dispersion method of
Linnik~\cite{BFI1986}, or the twisted fourth-moment method
of Deshouillers--Iwaniec.
\end{problem}

\begin{problem}[Computational validation]\label{prob:T4}
Compute $\mathcal{V}(N,Q)$ directly from prime-pair data for
$N \in \{10^8, 10^{10}, 10^{12}\}$, verifying that the
variance profile matches the prediction $\sim N^2/(Q(\log N)^4)$
and that the smooth-shift contribution is indeed subdominant.
\end{problem}

% ══════════════════════════════════════════════════════════

\subsection*{Acknowledgements}

The author thanks the anonymous reviewers of MT1 for identifying the
Rankin step as the critical heuristic and for recommending the
saddle-point approach as the natural rigorous replacement.

% ══════════════════════════════════════════════════════════
\begin{thebibliography}{99}

\bibitem{MT1}
S.~Durand,
\emph{Prime cluster bound via smooth number estimates:
unconditional upper bounds on close prime pairs},
MT1 preprint, Horizon Series, April 2026.

\bibitem{BFI1986}
E.~Bombieri, J.~B. Friedlander, and H.~Iwaniec,
\emph{Primes in arithmetic progressions to large moduli},
Acta Math. \textbf{156} (1986), 203--251.

\bibitem{Davenport2000}
H.~Davenport,
\emph{Multiplicative Number Theory},
3rd ed. (revised by H.\,L. Montgomery),
Springer, 2000.

\bibitem{Gallagher1976}
P.~X. Gallagher,
\emph{On the distribution of primes in short intervals},
Mathematika \textbf{23} (1976), 4--9.

\bibitem{GS2023}
A.~Granville and K.~Soundararajan,
\emph{Sums of singular series with large sets and the tail of the
distribution of primes},
Q. J. Math. \textbf{74} (2023), 1457--1475.

\bibitem{HT1986}
A.~Hildebrand and G.~Tenenbaum,
\emph{On integers free of large prime factors},
Trans. Amer. Math. Soc. \textbf{296} (1986), 265--290.

\bibitem{MV2007}
H.~L. Montgomery and R.~C. Vaughan,
\emph{Multiplicative Number Theory I: Classical Theory},
Cambridge University Press, 2007.

\bibitem{Polymath2014}
D.~H.~J. Polymath,
\emph{Variants of the Selberg sieve, and bounded intervals containing
many primes},
Res. Math. Sci. \textbf{1} (2014), Art.~12.

\bibitem{Tenenbaum2015}
G.~Tenenbaum,
\emph{Introduction to Analytic and Probabilistic Number Theory},
3rd ed., Amer. Math. Soc., 2015.

\bibitem{G36}
S.~Durand,
\emph{Horizon G36: Goldbach's conjecture reduces to a sharp prime
cluster bound},
Horizon Series, April 2026.

\bibitem{G37}
S.~Durand,
\emph{Horizon G37: The prime cluster bound --- canonical formulation,
Elliott--Halberstam conditional proof, and sieve-theoretic attack
strategy},
Horizon Series, April 2026.

\end{thebibliography}

\end{document}
